% Capacity transfer note, 2026-10-02. Only the earlier ordinary Sidon bound is Lean-checked.
\documentclass[10pt]{article}
\usepackage[hmargin=1in,vmargin=0.9in]{geometry}
\usepackage{amsmath,amssymb,amsthm}
\usepackage{booktabs}
\usepackage{microtype}
\usepackage[hidelinks]{hyperref}
\hypersetup{pdftitle={Interval capacity and second-order bounds for distinct-difference configurations},pdfauthor={Haoyu Chen},pdfsubject={Sonar sequences; weak Sidon sets; distinct differences}}
\numberwithin{equation}{section}
\newtheorem{theorem}{Theorem}[section]
\newtheorem{lemma}[theorem]{Lemma}
\newtheorem{proposition}[theorem]{Proposition}
\theoremstyle{remark}
\newtheorem{remark}[theorem]{Remark}
\newcommand{\R}{\mathbb R}
\newcommand{\Z}{\mathbb Z}
\newcommand{\ind}{\mathbf1}
\newcommand{\E}{\mathcal E}
\newcommand{\TV}{\mathrm{TV}}
\newcommand{\gam}{\frac{2\sqrt2}{3}}
\newcommand{\RefereeStatus}{Four isolated Claude (Anthropic) reports have been received: on cosine sonar; the one-dimensional transfers; the two-dimensional transfers and perturbation; and this whole paper, including the all-$N$ square supplement. The first three passed their respective mathematical claims, with requested scope and editorial repairs incorporated. The whole-paper verdict was PASS-WITH-REPAIRS: no mathematical error was found in a theorem, proposition or lemma, and its notational and wording repairs have been applied. Claude also supplied the original scout derivations. These reviews are cross-vendor relative to the GPT-6 Astra (OpenAI) proofs, but share the scout's vendor; no vendor-independent or human review is claimed.}
\makeatletter
\renewcommand\section{\@startsection{section}{1}{\z@}{-2.6ex \@plus -.8ex \@minus -.2ex}{1.3ex \@plus .2ex}{\normalfont\large\bfseries}}
\renewcommand\subsection{\@startsection{subsection}{2}{\z@}{-2.2ex \@plus -.6ex \@minus -.2ex}{1ex \@plus .2ex}{\normalfont\normalsize\bfseries}}
\makeatother
\title{Interval capacity and second-order bounds\\for distinct-difference configurations}
\author{Haoyu Chen\\Independent researcher}
\date{October 2, 2026}
\begin{document}
\maketitle
\begin{abstract}
We apply a signed interval-capacity certificate with boundary constant $2/3$
to several distinct-difference problems. A sonar sequence with $n$ rows and
$m$ columns satisfies
$m\le n+3(\pi^2/36)^{1/3}n^{2/3}+4n^{1/3}$ for $n\ge160^3$.
The argument preserves the exact column marginal in Cauchy--Schwarz.
We also give explicit bounds for weak Sidon sets, bounded difference
multiplicity, difference triangle sets, Manhattan configurations and
Sidon sets in integer boxes. All remainders and onsets are specified.
For a mass-one kernel $g$, the horizontal functional
$g(0)\int|t|g(t)\,dt$ has minimum $\pi^2/32$ among autocorrelations
$g=h*\widetilde h$ of nonnegative $h$ with $\int h=1$.
An explicit nonnegative positive-definite perturbation has a smaller
value, so that constant is not the infimum in the full admissible class;
the latter infimum remains open. The bounded-multiplicity, weak Sidon,
triangular sonar and difference-triangle bounds have been kernel-checked
in Lean~4 with Mathlib, as has the earlier ordinary Sidon bound; the
cosine sonar bound, the Manhattan and box bounds and the kernel
functional results are not formalised.
\end{abstract}
\noindent\textit{Keywords:} sonar sequences; Sidon sets; distinct differences;
interval capacity; signed measures; explicit bounds.

\section{Results and source-qualified comparisons}\label{sec:results}
We write $[N]=\{0,\ldots,N-1\}$; translation gives the same results for
$\{1,\ldots,N\}$. A strong Sidon set in $\Z^d$ has distinct unordered
pair sums, including repeated elements. Equivalently, every nonzero
ordered difference has at most one representation: from $a-b=a'-b'\ne0$
apply sum uniqueness to $a+b'=a'+b$, and conversely rearrange an equality
of sums. All roots are positive real roots. Integer parameters need not
be perfect powers. The bounds below are uniform over the stated sets.

A sonar sequence with $n$ rows and $m$ columns consists of
$y_0,\ldots,y_{m-1}\in[n]$, one point $(i,y_i)$ in each column, with
all nonzero ordered displacement vectors distinct. Repeated row values
are permitted; the differences are ordinary integer differences.
\begin{theorem}[Sonar]\label{thm:sonar}
For every such sequence,
\begin{align}
 m&\le n+2n^{2/3}+3n^{1/3} &&(n\ge48^3),\label{eq:sonar-ramp}\\
 m&\le n+3(\pi^2/36)^{1/3}n^{2/3}+4n^{1/3} &&(n\ge160^3).\label{eq:sonar-cos}
\end{align}
\end{theorem}
The best proved explicit published sonar bound located in the bounded
source audit is $m\le n+3.78n^{2/3}+4.76n^{1/3}+2$, for all $n$,
in Osorio--Ruiz--Trujillo--Urbano~\cite[Theorem~2.1]{ORTU14}.
Erd\H{o}s--Graham--Ruzsa--Taylor~\cite[Theorem~4]{EGRT92} state the
coefficient $5$; their printed proof gives $4$ with lower-order terms,
and their coefficient $3$ occurs in a remark without a proof.
Our comparison concerns the second-order coefficient and does not
assert pointwise superiority at every finite $n$.\footnote{The main
pages 104--108 of Chen--Kl\o ve~\cite{CK96}, the sonar part of
Robinson~\cite{Rob85}, and the Caicedo thesis were not available in the
source audit. The accessible Chen--Kl\o ve pages identify sonar as
$(h,p)=(2,1)$ and say their result generalises the EGRT bound; its exact
specialisation was not determined. These access gaps preclude an
exhaustive record or priority claim. Some sources interchange the
letters for rows and columns; our convention is fixed above.}

A weak Sidon set requires uniqueness only for unordered pairs of
\emph{distinct} elements; diagonal sums are excluded from its definition.
\begin{theorem}[Weak Sidon]\label{thm:weak}
For every integer $N\ge90^4$ and every weak Sidon $A\subseteq[N]$,
\begin{equation}\label{eq:weak}
 |A|<\sqrt N+\sqrt{8/3}\,N^{1/4}+2.
\end{equation}
\end{theorem}
The best proved coefficient of $N^{1/4}$ located in that audit is $\sqrt3-\eta$,
with $\eta\ge0.0089$, in a bound with $O(1)$ remainder, from
Balogh--F\"uredi--Roy~\cite[Theorem~5.1]{BFR21}. This is a citation to
arXiv version~2, not an assertion that the Monthly version contains that theorem.

\begin{theorem}[Bounded difference multiplicity]\label{thm:gthin}
Let $g,N\ge1$ be integers, $gN\ge120^4$, and $A\subseteq[N]$. If
$\#\{(a,b)\in A^2:a-b=d\}\le g$ for every nonzero integer $d$, then
\begin{equation}\label{eq:gthin}
 |A|<\sqrt{gN}+\gam(gN)^{1/4}+1.
\end{equation}
The same statement holds with $N$ replaced by the diameter $D$ of $A$
when $gD\ge120^4$.
\end{theorem}
This is the difference-multiplicity convention, not $B_2[g]$ sum
multiplicity. Balogh--F\"uredi--Roy~\cite[Section~6, (6.1)]{BFR21}
prove the second-order coefficient $1$ (of $(gN)^{1/4}$) with additive $1/2$.
Carter--Hunter--O'Bryant~\cite{CHO25} give a smaller, explicitly
$g$-dependent second-order coefficient $1-\varepsilon_g/2$; their journal abstract
provides $\varepsilon_g\ge0.0062g^{-4}$. The arXiv and journal bounds
are kept distinct in this comparison. Theorem~\ref{thm:gthin} is a
routine rescaling of the earlier scalar Sidon inequality.

An $(n,k)$ difference triangle set consists of $n$ rows of $k+1$
increasing integer marks, each row starting at zero, with all positive
within-row differences globally distinct. Its scope is the largest
mark, and $m(n,k)$ is the least possible scope. Cross-row differences
are not constrained.
\begin{theorem}[Difference triangles]\label{thm:dts}
For all integers $n\ge1$, $k\ge20365$, put
$\Phi(k)=(\sqrt{4k+8/9}-2\sqrt2/3)/2$. Then
\begin{align}
 m(n,k)&>n\Phi(k)^4,\label{eq:dts-inverse}\\
 m(n,k)&\ge n\left(k^2-\frac{4\sqrt2}{3}k^{3/2}
                      +\frac{16}{9}k-2\sqrt k\right).\label{eq:dts-expanded}
\end{align}
\end{theorem}
The located comparison is Kl\o ve's bound
$n(k^2-2k\sqrt k+(k+\sqrt k)/4)$, as quoted by
Chee--Colbourn~\cite[Theorem~1]{CC97} from Kl\o ve's 1988 Theorem~2.
This is a secondary-source comparison: the 1988 and 1989 primary
papers and Chen--Kl\o ve's 1991 lower-bound paper were not fully read.
Theorem~\ref{thm:dts} is another routine consequence of the scalar
inequality, summed over rows and inverted at a real parameter.

\begin{theorem}[Manhattan configurations]\label{thm:manhattan}
Let $A\subset\Z^2$ have unique nonzero ordered difference vectors,
and $|p_1-q_1|+|p_2-q_2|\le r$ for all $p,q\in A$.
For every real $r\ge160^3$,
\begin{equation}\label{eq:manhattan}
 |A|\le\frac r{\sqrt2}+(4/3)^{1/3}r^{2/3}+9r^{1/3}.
\end{equation}
\end{theorem}
The best bound located for this comparison is
$r/\sqrt2+3\,2^{-4/3}r^{2/3}+O(r^{1/3})$, in
Blackburn--Etzion--Martin--Paterson~\cite[Theorem~9, arXiv version~2]{BEMP10}.
The theorem numbering of the journal version was not verified.

\begin{theorem}[Integer boxes]\label{thm:boxes}
Let $d\ge2,N\ge1$ be integers, let $A\subseteq[N]^d$ be strong Sidon,
and set
\begin{equation}\label{eq:box-parameters}
 x=N^{d/(2d+2)},\qquad c_d=\frac{d+1}{2}(8/9)^{d/(d+1)}.
\end{equation}
If $x\ge\max(120,4d)$, then
\begin{equation}\label{eq:boxes}
 |A|\le N^{d/2}+c_dN^{d^2/(2d+2)}+2d^2N^{d(d-1)/(2d+2)}.
\end{equation}
For $d=2$, every integer $N\ge1$ also satisfies
\begin{equation}\label{eq:boxes-all}
 |A|\le N+(8/3)^{1/3}N^{2/3}+18N^{1/3}.
\end{equation}
\end{theorem}
For each fixed $d$ the final term in~\eqref{eq:boxes} is lower order;
no asymptotic uniformity in growing $d$ is asserted beyond its explicit
hypotheses. In dimension two the located coefficients of $N^{2/3}$ are $3/2$
asymptotically, attributed to Robinson~\cite{Rob85}, and $1.9$ with
remainder $1.6N^{1/3}+1$ for all $N$, attributed to Caicedo's 2016 thesis.
Both were read only as secondary statements in Trujillo~\cite[pp.~1038--1039]{Tru23}.
Bound~\eqref{eq:boxes-all} is smaller than that all-$N$ bound only for
$N\ge32433$; no pointwise superiority for all $N$ is asserted.
For $d\ge3$ the bounded audit did not establish the status of earlier
explicit constants. These statements are source comparisons, not a
claim of an exhaustive search.

The main additional inputs in this note are the exact column marginal
and the product use of a signed capacity certificate. The scalar
one-dimensional machinery comes from the ordinary Sidon note~\cite{Chen26}.
None of these bounds resolves the original conjecture on the
second-order term for ordinary Sidon sets.

\section{A self-contained signed capacity certificate}\label{sec:capacity}
We reproduce the analytic construction from~\cite{Chen26}. It supplies
unit potential on a closed interval, with a finite signed measure of
energy at most $L+2/3+200e^{-\alpha L}$, $\alpha=\log(4/3)$.
The ordinary Sidon consequence of that paper has been checked in Lean.
Apart from that earlier Sidon bound, no statement in this note is
claimed to be formalised in Lean.
\subsection{Kernel and signed energy}\label{sec:kernel}

Lebesgue measure on the real line is denoted by \(dt\). For an integrable
function \(v\) and a measure \(\mu\), write

\[
(v*\mu)(x)=\int v(x-t)\,d\mu(t)
\]

whenever the integral exists. A finite signed real measure means one with
finite total variation \(\|\mu\|_{\rm TV}=|\mu|(\mathbb R)\).

Set

\begin{equation}
h(t)=2(1-t)\mathbf1_{[0,1]}(t),\qquad
\widetilde h(t)=h(-t),\qquad f=h*\widetilde h.
\label{eq:s2-1}
\end{equation}

The specified value \(h(0)=2\) will be convenient in a pointwise identity.

\begin{lemma}[kernel and Cauchy inequality]\label{lem:2} The function \(f\) is even and

\begin{equation}
f(x)=
\begin{cases}
\displaystyle\frac43-2|x|+\frac23|x|^3,&|x|\le1,\\
0,&|x|>1.
\end{cases}
\label{eq:s2-2}
\end{equation}

It is nonnegative, nonincreasing on \([0,\infty)\), and satisfies

\begin{equation}
f(0)=\|h\|_2^2=\frac43,
\qquad \int_{\mathbb R}h(t)\,dt=1,
\qquad \int_{\mathbb R}f(t)\,dt=1.
\label{eq:s2-3}
\end{equation}

For finite signed real measures define

\[
E(\mu,\nu)=\iint f(x-y)\,d\mu(x)\,d\nu(y).
\]

Then

\begin{equation}
E(\mu,\nu)=\int_{\mathbb R}(h*\mu)(t)(h*\nu)(t)\,dt,
\qquad E(\mu,\mu)\ge0,
\label{eq:s2-4}
\end{equation}

and

\begin{equation}
|E(\mu,\nu)|^2\le E(\mu,\mu)E(\nu,\nu).
\label{eq:s2-5}
\end{equation}

\end{lemma}

\begin{proof} The autocorrelation is even by a change of variable. For
\(0\le x\le1\), its value is

\[
4\int_0^{1-x}(1-s-x)(1-s)\,ds
=\frac43-2x+\frac23x^3;
\]

for \(x>1\) the two supports have no overlap of positive length. This proves
\eqref{eq:s2-2}. Nonnegativity also follows directly from the integral of the two
nonnegative factors. On \([0,1]\) the derivative is \(-2+2x^2\le0\), and
outside this interval the function is zero. Integrating \(h\) and \(h^2\)
gives \eqref{eq:s2-3}; the mass of \(f\) is \((\int h)^2=1\) by Tonelli's theorem.

For any finite signed measure, Cauchy--Schwarz with respect to \(|\mu|\)
gives

\[
\int |h*\mu|^2\le
\|\mu\|_{\rm TV}
\int\!\int |h(t-x)|^2\,d|\mu|(x)\,dt
=\|\mu\|_{\rm TV}^2\|h\|_2^2.
\]

Thus \(h*\mu\in L^2\). Moreover,
\(\int |h(t-x)h(t-y)|\,dt\le\|h\|_2^2\), so the corresponding triple
integral against \(|\mu|\) and \(|\nu|\) is finite. Fubini's theorem is
therefore applicable. Its inner integral is \(f(x-y)\), which proves
\eqref{eq:s2-4}. Applying the usual Cauchy--Schwarz inequality in real \(L^2\) proves
\eqref{eq:s2-5}. This argument allows either measure to be signed. \end{proof}

\subsection{The half-line correction}\label{sec:halfline}

Let \(U\) be the probability measure with density
\(v_1=\mathbf1_{[0,1)}\), and let \(U^{*0}=\delta_0\). Define

\begin{equation}
g_+=\frac12\sum_{n=0}^{\infty}U^{*n}.
\label{eq:s3-1}
\end{equation}

The following construction proves everything needed about this measure;
no renewal theorem is used.

\begin{lemma}[local finiteness and exact potential]\label{lem:3} The measure \(g_+\) is
locally finite, supported on \([0,\infty)\), and

\begin{equation}
h*g_+=\mathbf1_{[0,\infty)},
\qquad (f*g_+)(x)=1\quad(x\ge0).
\label{eq:s3-2}
\end{equation}

\end{lemma}

\begin{proof} For \(n\ge1\), let \(v_n\) be the density of \(U^{*n}\).
Convolution gives representatives satisfying, for \(t\ge0\),

\begin{equation}
0\le v_n(t)\le\frac{t^{n-1}}{(n-1)!}.
\label{eq:s3-3}
\end{equation}

For \(n=1\) this is the bound \(v_1\le1\). If it holds at \(n\), then

\[
v_{n+1}(t)=\int_0^{\min(1,t)}v_n(t-s)\,ds
\le\int_0^t\frac{(t-s)^{n-1}}{(n-1)!}\,ds
=\frac{t^n}{n!}.
\]

For every fixed \(R\ge0\), it follows that
\(U^{*n}([0,R])\le R^n/n!\). The sum in \eqref{eq:s3-1} consequently has finite mass
on every bounded interval.

For a direct proof of the first identity in \eqref{eq:s3-2}, put
\(H=\mathbf1_{[0,\infty)}\). The definitions give the pointwise identity
\(h/2=H-U*H\). Thus each finite partial sum telescopes:

\begin{equation}
\frac12\sum_{n=0}^m h*U^{*n}
=H-U^{*(m+1)}*H.
\label{eq:s3-4}
\end{equation}

For \(t\ge0\), the last term on the right is between zero and
\(t^{m+1}/(m+1)!\), by \eqref{eq:s3-3}, and this tends to zero. For \(t<0\), both
sides vanish. The nonnegative partial sums on the left converge to
\(h*g_+\), proving the identity. This also fixes its value at \(t=0\).

Tonelli's theorem allows association of the nonnegative convolutions.
For \(x\ge0\),

\[
(f*g_+)(x)
=(\widetilde h*(h*g_+))(x)
=\int_{-1}^0\widetilde h(y)H(x-y)\,dy
=\int_{-1}^0\widetilde h(y)\,dy=1.
\]

All the integrals against \(g_+\) here are locally finite because their
kernels have compact support. \end{proof}

Write

\begin{equation}
g_+=\frac12\delta_0+\frac12u(t)\mathbf1_{[0,\infty)}(t)\,dt,
\qquad u(t)=\sum_{n=1}^{\infty}v_n(t)\quad(t\ge0).
\label{eq:s3-5}
\end{equation}

Use the right-continuous representative of \(u\). In particular,

\begin{equation}
u(0)=1,\quad u(t)=e^t\ (0\le t<1),\quad u(1)=e-1.
\label{eq:s3-6}
\end{equation}

The left limit at 1 is \(e\), and must not be substituted for \(u(1)\) in
the recurrence below. To justify these regularity claims, the series in
\eqref{eq:s3-5} converges uniformly on compact intervals by \eqref{eq:s3-3}. For \(n\ge2\),
the density \(v_n(t)=\int_{t-1}^t v_{n-1}(s)\,ds\), with densities set
to zero for \(s<0\), is continuous. Only \(v_1\) has a jump at 1. For
\(0\le t<1\), the induction proving \eqref{eq:s3-3} gives equality. At \(t=1\),
\(v_1(1)=0\) and \(v_n(1)=1/(n-1)!\) for \(n\ge2\). These observations
prove \eqref{eq:s3-6} and the stated representative convention.

\begin{lemma}[quantitative convergence of the density]\label{lem:4} For almost every
\(t\ge0\),

\begin{equation}
|u(t)-2|\le(e-1)(3/4)^{\lfloor t\rfloor}
<2(3/4)^{\lfloor t\rfloor}.
\label{eq:s3-7}
\end{equation}

\end{lemma}

\begin{proof} Summing the convolution recurrence for the densities yields
\(u=v_1+U*u\). Thus

\begin{equation}
u(t)=\int_{t-1}^t u(s)\,ds\qquad(t\ge1).
\label{eq:s3-8}
\end{equation}

This includes \(t=1\) with the right-continuous value in \eqref{eq:s3-6}. On
\((1,\infty)\), this identity makes \(u\) locally absolutely continuous,
and \(u'(t)=u(t)-u(t-1)\) almost everywhere. Fix an integer \(n\ge1\).
Solving this first-order equation on \([n,n+1]\) gives

\[
u(n+x)=e^xu(n)-\int_0^x e^{x-y}u(n-1+y)\,dy
\qquad(0\le x\le1).
\]

Since \(u(n)=\int_0^1u(n-1+y)\,dy\), this becomes

\begin{equation}
u(n+x)=\int_0^1K_x(y)u(n-1+y)\,dy,
\qquad
K_x(y)=e^x-\mathbf1_{\{y\le x\}}e^{x-y}.
\label{eq:s3-9}
\end{equation}

For \(0\le x,y\le1\), the kernel is nonnegative, and

\[
\int_0^1 K_x(y)\,dy=e^x-\int_0^x e^{x-y}\,dy=1.
\]

For \(y\in[1/2,1]\), it also satisfies \(K_x(y)\ge1/2\). Indeed, if
\(y>x\), its value is \(e^x\ge1\). If \(y\le x\), its value is
\(e^{x-y}(e^y-1)\ge e^{1/2}-1>1/2\), using the exponential series.

Let \([m_n,M_n]\) be the smallest closed interval containing the essential
range of \(u\) on \([n,n+1)\). Here essential extrema ignore sets of
Lebesgue measure zero. Initially \([m_0,M_0]=[1,e]\). Equation \eqref{eq:s3-9} first
implies that these closed intervals are nested. For the width estimate,
every \(K_x(y)\,dy\) contains the common measure

\[
w(y)\,dy=\frac12\mathbf1_{[1/2,1]}(y)\,dy
\]

of mass \(1/4\). Put
\(c_n=\int_0^1w(y)u(n-1+y)\,dy\). The remaining nonnegative measure
has mass \(3/4\), so every value in \eqref{eq:s3-9} lies between
\(c_n+(3/4)m_{n-1}\) and \(c_n+(3/4)M_{n-1}\). Consequently,

\[
[m_n,M_n]\subseteq[m_{n-1},M_{n-1}],
\qquad M_n-m_n\le\frac34(M_{n-1}-m_{n-1}).
\]

The nested closed intervals therefore meet in a single real number \(c\),
and

\begin{equation}
|u(t)-c|\le(e-1)(3/4)^{\lfloor t\rfloor}
\label{eq:s3-10}
\end{equation}

almost everywhere. This also provides a uniform essential bound and
essential uniform convergence to \(c\) as \(t\to\infty\).

For \(t>1\), the atom in \eqref{eq:s3-5} contributes nothing to \(h*g_+(t)\), so
Lemma~\ref{lem:3} gives

\[
1=\frac12\int_0^1 h(s)u(t-s)\,ds.
\]

Letting \(t\to\infty\) in this integral, \eqref{eq:s3-10} and \(\int h=1\) imply
\(1=c/2\). Thus \(c=2\). Finally, \(e<3\): in its series
\(e=2+\sum_{j\ge2}1/j!\), the tail is strictly smaller than
\(\sum_{j\ge2}2^{-(j-1)}=1\). This proves \eqref{eq:s3-7}. \end{proof}

\begin{lemma}[size, tail, and mass of the correction]\label{lem:5} Define the signed
measure supported on \([0,\infty)\)

\begin{equation}
q=g_+-\mathbf1_{[0,\infty)}(t)\,dt
=\frac12\delta_0+r(t)\mathbf1_{[0,\infty)}(t)\,dt,
\qquad r(t)=\frac12u(t)-1.
\label{eq:s3-11}
\end{equation}

It is a finite signed measure, and, with \(\alpha=\log(4/3)\),

\begin{equation}
|r(t)|\le(3/4)^{\lfloor t\rfloor}\quad\hbox{almost everywhere},
\qquad \|q\|_{\rm TV}\le\frac92,
\label{eq:s3-12}
\end{equation}

\begin{equation}
|q|((L,\infty))\le6e^{-\alpha L}\qquad(L\ge1),
\label{eq:s3-13}
\end{equation}

and

\begin{equation}
q(\mathbb R)=\frac13.
\label{eq:s3-14}
\end{equation}

\end{lemma}

\begin{proof} The first bound follows from Lemma~\ref{lem:4}. Integration on unit intervals
then gives

\[
\|q\|_{\rm TV}\le\frac12+\sum_{n=0}^\infty(3/4)^n=\frac92.
\]

For \(L\ge1\), the atom at zero is outside \((L,\infty)\), and

\[
|q|((L,\infty))
\le\sum_{n=\lfloor L\rfloor}^\infty(3/4)^n
=4(3/4)^{\lfloor L\rfloor}
\le\frac{16}{3}e^{-\alpha L}
\le6e^{-\alpha L}.
\]

It remains to compute the signed mass, not just its absolute bound. For
\(s>0\), Tonelli's theorem and the product rule for Laplace transforms of
nonnegative convolutions give

\begin{equation}
\begin{aligned}
\int e^{-st}\,dg_+(t)
&=\frac12\sum_{n=0}^\infty
\left(\int_0^1e^{-st}\,dt\right)^n\\
&=\frac1{2\left(1-(1-e^{-s})/s\right)}.
\end{aligned}
\label{eq:s3-15}
\end{equation}

The geometric series converges since \(0<\int_0^1e^{-st}\,dt<1\).
Taylor's formula for the exponential, as \(s\downarrow0\), gives

\[
\phi(s):=1-\frac{1-e^{-s}}s
=\frac s2-\frac{s^2}6+O(s^3).
\]

This use of \(O(s^3)\) concerns only a limit: for \(0<s\le1\), Taylor's
remainder after the cubic term of \(e^{-s}\) has absolute value at most
\(s^4/24\), which supplies the displayed bound directly. Subtracting the
Laplace transform \(1/s\) of half-line Lebesgue measure gives

\[
\int e^{-st}\,dq(t)
=\frac1{2\phi(s)}-\frac1s
=\frac{s-2\phi(s)}{2s\phi(s)}
\longrightarrow\frac13.
\]

Because \(q\) has finite total variation and is supported on the
nonnegative half-line, \(|e^{-st}|\le1\) there. Dominated convergence for
the finite measure \(|q|\) shows that the limit is \(q(\mathbb R)\), proving
\eqref{eq:s3-14}. \end{proof}

\subsection{Joining two half-lines}\label{sec:finite}

\begin{lemma}[finite-interval energy bound]\label{lem:6} Let \(L\ge1\), and let \(\mu\)
be any finite positive measure supported on the closed interval \([0,L]\).
If \(k=\mu(\mathbb R)\), then

\begin{equation}
E(\mu,\mu)\ge
\frac{k^2}{L+\frac23+200e^{-\alpha L}},
\qquad \alpha=\log(4/3).
\label{eq:s4-1}
\end{equation}

\end{lemma}

\begin{proof} Let \(q_L\) be the pushforward of \(q\) by the reflection
\(t\mapsto L-t\), and set

\begin{equation}
\nu_L=\mathbf1_{[0,L]}(t)\,dt+q+q_L.
\label{eq:s4-2}
\end{equation}

This is a finite signed real measure, with

\begin{equation}
\nu_L(\mathbb R)=L+\frac23.
\label{eq:s4-3}
\end{equation}

Let \(g_-^L\) denote the same reflection of \(g_+\). As locally finite
measures,

\begin{equation}
\nu_L=g_++g_-^L-dt.
\label{eq:s4-4}
\end{equation}

Indeed, the two half-line Lebesgue measures sum to full-line Lebesgue
measure plus Lebesgue measure on \([0,L]\); endpoints do not change those
Lebesgue measures. The atoms of \(q\) at 0 and of \(q_L\) at \(L\) remain
included in \eqref{eq:s4-2}. Since \(f\) is compactly supported, \eqref{eq:s4-4} may be
convolved locally with \(f\). Lemma~\ref{lem:3} and evenness of \(f\) give, for
\(0\le x\le L\),

\begin{equation}
V_L(x):=(f*\nu_L)(x)=1+1-\int f=1.
\label{eq:s4-5}
\end{equation}

This holds also at both endpoints. For every real \(x\), nonnegativity
and mass one of \(f\), together with Lemma~\ref{lem:5}, imply

\begin{equation}
|V_L(x)|\le1+\|f\|_\infty
(\|q\|_{\rm TV}+\|q_L\|_{\rm TV})
\le1+\frac43\cdot9=13.
\label{eq:s4-6}
\end{equation}

Outside \([0,L]\), the only measure in \eqref{eq:s4-2} on \((L,\infty)\) is the
tail of \(q\); on \(( -\infty,0)\) it is the reflection of the same tail.
In particular there are no excluded endpoint atoms. Thus \eqref{eq:s3-13} gives

\[
|\nu_L|(\mathbb R\setminus[0,L])\le12e^{-\alpha L}.
\]

By \eqref{eq:s4-3}, \eqref{eq:s4-5}, and \eqref{eq:s4-6},

\begin{equation}
\begin{aligned}
\left|E(\nu_L,\nu_L)-\left(L+\frac23\right)\right|
&=\left|\int_{\mathbb R\setminus[0,L]}(V_L(x)-1)\,d\nu_L(x)\right|\\
&\le14\cdot12e^{-\alpha L}
=168e^{-\alpha L}
\le200e^{-\alpha L}.
\end{aligned}
\label{eq:s4-7}
\end{equation}

Also \(E(\mu,\nu_L)=\int V_L\,d\mu=k\). Lemma~\ref{lem:2} now gives

\[
k^2\le E(\mu,\mu)E(\nu_L,\nu_L)
\le E(\mu,\mu)\left(L+\frac23+200e^{-\alpha L}\right).
\]

The denominator is positive and \(E(\mu,\mu)\ge0\), proving \eqref{eq:s4-1}.
There is no assumption that \(\nu_L\) is positive. \end{proof}

\section{Sonar and the exact column marginal}\label{sec:sonar}
Write $a_2=4/3$, $b=2/3$, $\alpha=\log(4/3)$ and
$C(L)=L+b+200e^{-\alpha L}$. Let $g$ be an even nonnegative
positive-definite horizontal kernel, $\mu=\sum_i\delta_{(i,y_i)}$,
$\lambda=\sum_{i=0}^{m-1}\delta_i$, and let $\rho$ be the pushforward
of $\nu_{n/U}$ by $t\mapsto Ut$, without multiplying its mass.
For a bounded kernel $K$ on $\R^j$ and finite signed measures write
$E_K(\mu,\nu)=\iint K(x-y)\,d\mu(x)\,d\nu(y)$; thus $E_f=E$
of Lemma~\ref{lem:2}.
For $G(x,y)=g(x/T)f(y/U)$ set
\begin{equation}\label{eq:marginal}
 \Lambda=\sum_{i,j=0}^{m-1}g((i-j)/T),\qquad
 E_G(\mu,\lambda\otimes\rho)=\Lambda,\qquad
 E_G(\lambda\otimes\rho,\lambda\otimes\rho)
       =\Lambda E_f(\nu_{n/U},\nu_{n/U}).
\end{equation}
When $m>0$, $g(0)>0$ ensures $\Lambda>0$. Signed Cauchy--Schwarz and
Lemma~\ref{lem:6} give
\begin{equation}\label{eq:marginal-bound}
 \Lambda\le C(n/U)E_G(\mu,\mu)\qquad(0<U\le n).
\end{equation}
For the autocorrelations below this is an $L^2$ Gram inequality: the
scaled factors are $T^{-1/2}h_g(\cdot/T)$ and
$U^{-1/2}h(\cdot/U)$, where $g=h_g*\widetilde h_g$; here
$h_g=\ind_{[0,1]}$ for the triangle and $h_g=h_c$ for the cosine
kernel below. Bounded kernels and finite total variations
justify the product integrals. Kernel nonnegativity, separately,
allows completion of missing differences. The exact column marginal
makes~\eqref{eq:marginal} possible. A prescribed marginal in another
product problem gives the same argument; an unrestricted projection
cannot be substituted for a prescribed one.

\subsection{The triangular horizontal kernel}
Take $g(t)=(1-|t|)_+$ and $T=V$ for an integer $1\le V\le m$.
Exact summation gives
\begin{equation}\label{eq:triangle-sandwich}
 \Lambda=mV-\frac{V^2-1}{3},\qquad
 E_G(\mu,\mu)\le\frac43m+(V-1)(U+4/3).
\end{equation}
Indeed $\sum_{d\ne0}g(d/V)=V-1$ and
$\sum_{e\in\Z}f(e/U)\le U+4/3$. Every off-diagonal sonar difference
has nonzero first coordinate and occurs once. The formula for
$\Lambda$ is $m+2\sum_{d=1}^{V-1}(m-d)(1-d/V)$.

Put $x=n^{1/3}\ge48$, $U=x^2$ and $V=\lceil2x^2\rceil\le n$.
If $m<V$ the result is immediate; otherwise the preceding inequalities apply.
As $\alpha\ge1/4$, $x^2e^{-\alpha x}$ decreases for $x\ge8$.
The exact comparisons $(3/4)^8<1/9$ and $200\cdot48^2<9^6$ imply
$\varepsilon=200e^{-\alpha x}<x^{-2}$. With $C=x+b+\varepsilon$,
solving the sandwich linearly gives
\begin{equation}\label{eq:triangle-linear}
 m(1-a_2C/V)\le V/3-1/(3V)+C(1-1/V)(x^2+a_2).
\end{equation}
Set
\[
 q=\frac2{3x}+\frac4{9x^2}+\frac2{3x^4},\qquad
 B=x^3+\frac43x^2+\frac56x+\frac{20}{9}+\frac1{4x}+\frac4{3x^2}.
\]
Then $a_2C/V\le q<1/2$, and the right side of
\eqref{eq:triangle-linear} is at most $B$.
For the latter use $V\le2x^2+1$, $C\le x+b+x^{-2}$, $C\ge x$, and
$x^3/(2x^2+1)\ge x/2-1/(4x)$. With $Y=x^3+2x^2+3x$,
\[
 36x^3\{Y(1-q)-B\}
 =14x^4-184x^3-81x^2-96x-72
 =14t^4+712t^3+12591t^2+85376t+141496,
 \quad t=x-16.
\]
It is positive for $x\ge48$. Division by $1-q>0$ proves
\eqref{eq:sonar-ramp} in Theorem~\ref{thm:sonar}, without a preliminary
asymptotic bound on $m$.

\subsection{The cosine autocorrelation}
Let $h_c(t)=(\pi/2)\sin(\pi t)\ind_{[0,1]}(t)$ and $g=h_c*\widetilde h_c$.
For $0\le s\le1$, direct integration gives
\begin{equation}\label{eq:cosine-kernel}
 g(s)=\frac{\pi^2}{8}(1-s)\cos(\pi s)+\frac\pi8\sin(\pi s),\qquad
 a=g(0)=\frac{\pi^2}{8},\qquad M=\int|s|g(s)\,ds=\frac14.
\end{equation}
Extend evenly and by zero outside $[-1,1]$. The mass is one and
$g'(s)=-(\pi^3/8)(1-s)\sin(\pi s)\le0$ on $[0,1]$.
The moment also follows from the cumulative distribution
$F(s)=(1-\cos(\pi s))/2$: $M=2\int F(1-F)=1/4$.

For any even decreasing nonnegative mass-one kernel $w$,
\[
 T-w(0)\le\sum_{d\in\Z}w(d/T)\le T+w(0),\qquad
 \sum_{d\ne0}w(d/T)\le T.
\]
For $p(s)=|s|g(s)$, integration on $(0,\infty)$ gives
$\int|p'|\le\int g+\int s|g'|=1$; hence $\operatorname{Var}(p)\le2$.
Comparing each sample to its adjacent cell integral yields
$|\sum_d p(d/T)-TM|\le2$. For $T\le m$ these estimates give
\begin{align}
 \Lambda&\ge m(T-a)-MT^2-2T,\label{eq:cos-lower}\\
 E_G(\mu,\mu)&\le aa_2m+T(U+a_2).\label{eq:cos-upper}
\end{align}
Thus the fully finite, linear inequality is
\begin{equation}\label{eq:cos-linear}
 m[1-a/T-aa_2C(n/U)/T]\le MT+2+C(n/U)(U+a_2).
\end{equation}

Put $x=n^{1/3}$, $v=(\pi^2/36)^{1/3}$,
$T=4vx^2$ and $U=(3/2)vx^2$.
The elementary bounds $3<\pi<22/7$ give $3/5<v<2/3$ and
$a=(9/2)v^3$. If $m<n$ there is nothing to prove; otherwise
$x\ge160$ ensures $T\le n\le m$, $U\le n$.
As $\alpha=\int_1^{4/3}dt/t\ge1/4$,
\begin{equation}\label{eq:cos-tail}
 \varepsilon=200e^{-\alpha n/U}\le200e^{-x/4}\le x^{-4}
 \qquad(x\ge160).
\end{equation}
Indeed $x^4e^{-x/4}$ decreases for $x\ge16$ and
$200\cdot160^4=131072000000<2^{40}<e^{40}$.
Set $d_0=(17/8)v^2$. The denominator and right side in
\eqref{eq:cos-linear} are $1-u$ and $B_c$, where
\[
 u=\frac vx+\frac{d_0+(3/2)v^2\varepsilon}{x^2},\qquad
 B_c=x^3+2vx^2+\frac8{9v}x+\frac{26}{9}
                  +\varepsilon\left(\frac32vx^2+\frac43\right).
\]
We have $0<u<1/2$. For $Y=x^3+3vx^2+4x$, exact expansion gives
\begin{align*}
 Y(1-u)-B_c
 &=\left(4-\frac{41}{8}v^2-\frac8{9v}\right)x
    -4v-3vd_0-4d_0/x-26/9\\
 &\quad-\varepsilon\left(\frac32v^2x+\frac92v^3+
                    \frac{6v^2}{x}+\frac32vx^2+\frac43\right).
\end{align*}
The coefficient of $x$ is at least $13/54$,
$4v+3vd_0+26/9\le67/9$, and $4d_0/x\le1$.
The final error is at most
$x^{-2}+(2/3)x^{-3}+(8/3)x^{-4}+(8/3)x^{-5}<1$.
Consequently $Y(1-u)-B_c\ge13x/54-85/9>0$.
This proves~\eqref{eq:sonar-cos} for every real $x\ge160$.
Finally $3v<2$ follows from $(22/7)^2<32/3$, so the coefficient in
\eqref{eq:sonar-cos} is smaller than that in~\eqref{eq:sonar-ramp}.

\subsection{The first-moment functional and its precise scope}
\begin{proposition}\label{prop:factor}
For $h\ge0$ in $L^1(\R)\cap L^2(\R)$ with $\int h=1$, let
$g=h*\widetilde h$, $a=g(0)$ and $M=\int|t|g(t)\,dt$. Then
\begin{equation}\label{eq:factor-optimum}
 aM\ge\frac{\pi^2}{32}.
\end{equation}
Equality is attained precisely by translates and dilates of the sine
density $h_c$ above.
\end{proposition}
\begin{proof}
The case $M=\infty$ is immediate. With $F(x)=\int_{-\infty}^xh(t)\,dt$,
Tonelli applied to thresholds separating two points gives
$M=2\int F(1-F)$. Cauchy--Schwarz and the absolutely continuous
change of variables $u=F(x)$ give
\[
 aM=2\left(\int h^2\right)\left(\int F(1-F)\right)
 \ge2\left(\int h\sqrt{F(1-F)}\right)^2
 =2\left(\int_0^1\sqrt{u(1-u)}\,du\right)^2=\frac{\pi^2}{32}.
\]
The change of variables remains valid where $h$ vanishes.
Equality requires $h=c\sqrt{F(1-F)}$ almost everywhere.
On the interval where $0<F<1$, integration of this differential
equation gives the sine density on one interval; outside it $h=0$
almost everywhere. Translation and dilation are the only freedoms.
\end{proof}

After solving~\eqref{eq:cos-linear} for $m$, the excess over $n$ is
$MT+bU+aa_2n^2/(TU)$ plus lower-order terms. Its AM--GM minimum is
$3[(aM)(a_2b)]^{1/3}n^{2/3}$. For the fixed vertical ramp,
$a_2b=8/9$, so~\eqref{eq:factor-optimum} shows that, with the vertical
ramp fixed, no horizontal kernel that is the autocorrelation of a
nonnegative unit-mass factor makes the AM--GM minimum of this leading
expression smaller than $3(\pi^2/36)^{1/3}n^{2/3}$.
This does not concern non-product two-dimensional
kernels or inequalities using additional information.

For completeness, the separate one-dimensional theorem of~\cite{Chen26}
says that if $w\in C_0(\R)\cap L^1(\R)$ is even, nonnegative, of
mass one and $w(0)>0$, and $C_w(L)$ is the reciprocal minimum energy
of positive probability measures on $[0,L]$, then
$\liminf_{L\to\infty}(C_w(L)-L)\ge8/(9w(0))$.
Thus a finite intercept $C_w(L)=L+b_w+o(1)$ satisfies
$w(0)b_w\ge8/9$. Varying the vertical kernel preserves the preceding
template conclusion only under these hypotheses and the sampling
and certificate asymptotics needed to justify the template.

\begin{proposition}\label{prop:full-lower}
Every continuous, even, nonnegative positive-definite integrable
kernel $g$ with mass one satisfies
\begin{equation}\label{eq:full-lower}
 g(0)\int|t|g(t)\,dt\ge\frac14+\frac1{18\pi^2}.
\end{equation}
\end{proposition}
\begin{proof}
Put $a=g(0)>0$ and $q(t)=g(t/a)/a$. Then $q(0)=\int q=1$ and
$\int|t|q=a\int|t|g$. Positive definiteness gives $0\le q\le1$
and $\widehat q\ge0$, with Fourier phase $e^{-2\pi i\xi t}$.
The latter follows either from the positive-definite quadratic form
and truncated exponentials, or from its Fourier representation.
For $B=\ind_{[-1/2,1/2]}$, $\widehat B(3/2)=-2/(3\pi)$, whence
$\|q-B\|_1\ge2/(3\pi)$. Let
$\eta=\int_{|t|>1/2}q=\int_{|t|\le1/2}(1-q)\ge1/(3\pi)$.
Then
\[
 \int|t|q-\frac14
 =\int_{|t|>1/2}(|t|-1/2)q
    +\int_{|t|\le1/2}(1/2-|t|)(1-q)
 \ge\frac{\eta^2}{2}.
\]
Each integral is at least $\eta^2/4$: for a density at most one,
the distance distribution from two boundary points has at most
$2s$ mass within distance $s$, so integrate $(\eta-2s)_+$.
The inner strips fit because $\eta\le1$.
\end{proof}

\subsection{An admissible perturbation below the restricted optimum}
\begin{proposition}\label{prop:perturb}
The constant $\pi^2/32$ is not the infimum over the full class in
Proposition~\ref{prop:full-lower}. In fact, set
\begin{equation}\label{eq:perturb-def}
 \delta=\frac1{16},\quad \theta=\frac1{10^7},\quad
 b_\delta(t)=\delta^{-1}g(t/\delta),\quad
 g_\theta(t)=g(t)+\theta[b_\delta(t-2)+b_\delta(t+2)-2b_\delta(t)],
\end{equation}
where $g$ is the cosine kernel~\eqref{eq:cosine-kernel}. This kernel
is continuous, even, nonnegative, positive definite, and of mass one,
and, with $R=1-(129/8)\theta-508\theta^2$,
\begin{equation}\label{eq:perturb-functional}
 g_\theta(0)\int|t|g_\theta(t)\,dt=\frac{\pi^2}{32}R<\frac{\pi^2}{32}.
\end{equation}
Consequently the full-class infimum lies in
$[1/4+1/(18\pi^2),(\pi^2/32)R]$; its exact value remains open.
\end{proposition}
\begin{proof}
Write $a_0=\pi^2/8$. The two positive tails are disjoint from
$[-1,1]$. The only negative correction is in $|t|\le\delta$, where
$g(t)\ge a_0(1-\delta)\cos(\pi\delta)\ge465a_0/512>a_0/2$,
using $\pi\delta<1/4$ and $\cos u\ge1-u^2/2$.
The subtraction is at most $32\theta a_0<a_0/2$.
Thus $g_\theta\ge0$, and its mass remains one. Direct moments give
\[
 a=g_\theta(0)=a_0(1-32\theta),\qquad
 M=\int|t|g_\theta(t)\,dt=\frac14+\frac{127}{32}\theta,
\]
which proves~\eqref{eq:perturb-functional} and $0<R<1$.

The Fourier transforms are
\[
 \widehat h_c(\xi)=\frac{e^{-\pi i\xi}\cos(\pi\xi)}{1-4\xi^2},\qquad
 \widehat g(\xi)=\frac{\cos^2(\pi\xi)}{(1-4\xi^2)^2},\qquad
 \widehat g_\theta(\xi)=\widehat g(\xi)
                -4\theta\sin^2(2\pi\xi)\widehat g(\delta\xi),
\]
with removable values at $\xi=\pm1/2$. The distributional second
derivative of the zero extension of $h_c$ has total variation $2\pi^2$:
$\pi^2$ from its endpoint derivative jumps and $\pi^2$ from the interior.
Hence $|\widehat h_c(\xi)|\le\min(1,1/(2\xi^2))$.
Where $\widehat g\ne0$, the ratio of the subtracted term before
$\theta$ to $\widehat g$ is
\[
 16\sin^2(\pi\xi)(1-4\xi^2)^2\widehat g(\delta\xi)
 \le16(1+2/\delta^2)^2=4210704.
\]
To obtain the bound, split at $\xi^2=1/(2\delta^2)$ and use respectively
$\widehat g\le1$ and $\widehat g(\delta\xi)\le1/(4\delta^4\xi^4)$.
Continuity handles the zeros and removable points. Since
$4210704\theta<1$, the transform is nonnegative. Its integrable
fourth-power decay justifies Fourier inversion and proves positive
definiteness. By Proposition~\ref{prop:factor}, this example cannot
have a nonnegative unit-mass factor in $L^1\cap L^2$.
\end{proof}

A secondary consequence, with a larger explicit remainder, is
\begin{equation}\label{eq:perturb-sonar}
 m\le n+3[(\pi^2/36)R]^{1/3}n^{2/3}+8n^{1/3},\qquad n\ge160^3.
\end{equation}
Here are the error details. For the perturbed kernel,
$1\le a\le4/3$, $1/4\le M\le1/3$, its support lies in $[-33/16,33/16]$,
and
\[
 \operatorname{Var}(g_\theta)\le2a_0+128\theta a_0<3,\qquad
 \operatorname{Var}(|t|g_\theta)\le2+\theta(6+132a_0)<3.
\]
For the second estimate, each shifted $|t|b_\delta(t\mp2)$ has
variation at most $1+(2+\delta)2a_0/\delta$, and the central term
$2|t|b_\delta(t)$ has variation at most $4$, twice the variation of
$|s|g(s)$.
Adjacent-cell comparisons give
$|\sum_d g_\theta(d/T)-T|\le3$ and
$\sum_d|d|g_\theta(d/T)\le MT^2+3T$.
The nonzero horizontal sum is at most $T+3-a\le T+2$.
The same marginal argument, now justified by the nonnegative integrable
Fourier transform for signed measures, yields
\begin{equation}\label{eq:perturb-linear}
 m[1-3/T-(4a/3)C/T]\le MT+3+C(1+2/T)(U+4/3),
 \quad C=n/U+2/3+\varepsilon,
\end{equation}
provided $(33/16)T\le m$, $U\le n$.
Put $x=n^{1/3}$, $v=[(\pi^2/36)R]^{1/3}$,
$t_1=v/M$, $t_2=3v/2$, $T=t_1x^2$, $U=t_2x^2$.
We have $3/5<v<2/3$, $9/5<t_1<8/3$, $9/10<t_2<1$;
these follow from $R>9/10$ and the preceding rational bounds.
If $m\ge n$ and $x\ge160$, the support condition holds since
$(33/16)t_1x^2\le(11/2)x^2\le n$. The tail again obeys
$\varepsilon\le x^{-4}$.
In~\eqref{eq:perturb-linear} the denominator is $1-u$, where
\[
 u=v/x+(d_1+e_1\varepsilon)/x^2,\quad
 d_1=\frac{3+8a/9}{t_1}\le\frac{565}{243}<3,\quad
 e_1=\frac{4a/3}{t_1}\le\frac{80}{81}<1.
\]
It is positive and $u<1/2$. The right side is at most
$x^3+2vx^2+(70/27)x+8$: the $x$ terms have coefficients
$4/(3t_2)+2/t_1\le70/27$;
the constants $3+8/9+2bt_2/t_1$ are at most $125/27$;
and the remaining terms
$\varepsilon t_2x^2+(4/3)\varepsilon
+2(4/3)/(t_1t_2x)+2(4/3)(b+\varepsilon)/(t_1x^2)
+2\varepsilon t_2/t_1$
are together less than $3$ at $x\ge160$.
For $Y=x^3+3vx^2+8x$, substitution of
$v\le2/3$, $d_1\le3$, $e_1\varepsilon\le x^{-4}$ gives
\[
 Y(1-u)-[x^3+2vx^2+(70/27)x+8]\ge\frac{29}{27}x-22>0.
\]
This proves~\eqref{eq:perturb-sonar}; the case $m<n$ is immediate.
The exact arithmetic bounds, rather than a fitted numerical optimum,
are used throughout. The smaller coefficient in
\eqref{eq:perturb-sonar} does not imply a smaller bound at every finite $n$.

\section{One-dimensional corollaries}\label{sec:one-dimensional}
For the ramp, monotonicity gives $2\sum_{d\ge1}f(d/T)\le T$
for every real $T>0$. We first record the real-parameter scalar step
so that subsequent substitutions do not assume integrality.
\begin{lemma}\label{lem:scalar}
Let $X\ge120^4$, $k\ge0$ be real, put $x=X^{1/4}$,
$S=\sqrt2x^3$, and suppose
\begin{equation}\label{eq:scalar}
 k^2\le\left(X/S+2/3+200e^{-\alpha X/S}\right)(S+4k/3).
\end{equation}
Then $k<x^2+(2\sqrt2/3)x+1$.
\end{lemma}
\begin{proof}
Write $\gamma=2\sqrt2/3$ and $\varepsilon=200e^{-\alpha x/\sqrt2}$.
The scale obeys $0<S\le X$, $X/S\ge1$. Since
$\alpha/\sqrt2\ge1/6$ and $xe^{-x/6}$ decreases for $x\ge6$,
\[
 \varepsilon x\le24000e^{-20}<24000/2^{20}<1/32.
\]
Expansion of~\eqref{eq:scalar} gives $P_\varepsilon(k)\le0$, where
\[
 P_\varepsilon(z)=z^2-(\gamma x+8/9+4\varepsilon/3)z
                   -x^4-\gamma x^3-\sqrt2\varepsilon x^3.
\]
At $y=x^2+\gamma x+1$,
\[
 P_0(y)=\frac{10}{9}x^2+\frac\gamma9x+\frac19\ge x^2,
 \qquad
 \varepsilon(4y/3+\sqrt2x^3)<x/8+3x^2/64\le11x^2/64.
\]
Here $y\le3x^2$. Thus $P_\varepsilon(y)>0$. Its constant term is
negative, so this monic quadratic has exactly one positive root,
which lies between $k$ and $y$.
\end{proof}

\subsection{Weak Sidon sets}
For $d>0$ let $r(d)=\#\{(x,y)\in A^2:x-y=d\}$. If two distinct
lower endpoints $p<q$ represent $d$, then
$(p+d)+q=(q+d)+p$. Unless $q=p+d$, both pairs consist of distinct
elements, are different, and contradict weak Sidonicity.
Thus repeated differences are exactly adjacent edges of a three-term
progression. Three lower endpoints are impossible, so $r(d)\le2$.
Distinct repeated differences have distinct middle points: otherwise
two distinct off-diagonal pairs have the same sum twice that point.
These middle points are neither the minimum nor the maximum of $A$.
Consequently, with $k=|A|\ge2$,
\begin{equation}\label{eq:weak-energy}
 E_f(\mu,\mu)\le\frac43k+T+\frac83(k-2)
                 =T+4k-16/3\le T+4k,
 \qquad \mu=\sum_{a\in A}\delta_{a/T}.
\end{equation}
For $k=0,1$ the last bound holds directly. The diagonal energy remains
$4k/3$ even though diagonal sums are excluded from weak Sidonicity.
Lemma~\ref{lem:6} therefore gives
\[
 k^2\le(N/T+2/3+200e^{-\alpha N/T})(T+4k).
\]
Put $x=N^{1/4}\ge90$, $T=\sqrt6x^3$ and $\beta=2\sqrt6/3=\sqrt{8/3}$.
Then $T\le N$ and, writing $\varepsilon=200e^{-\alpha x/\sqrt6}$,
\begin{equation}\label{eq:weak-poly}
 P_\varepsilon(k)\le0,\qquad
 P_\varepsilon(z)=z^2-(\beta x+\beta^2+4\varepsilon)z
                   -x^4-\beta x^3-\sqrt6\varepsilon x^3.
\end{equation}
The derivative of $\log u-2(u-1)/(u+1)$ is
$(u-1)^2/[u(u+1)^2]\ge0$, so $\alpha\ge2/7$.
As $\sqrt6<49/20$, $\alpha/\sqrt6\ge40/343$ and monotonicity gives
$\varepsilon x\le18000e^{-3600/343}<1/2$.
The strict last comparison has the exact rational certificate
\[
 (1957/720)^{10}\sum_{j=0}^4\frac{(170/343)^j}{j!}>36000,
 \qquad 3600/343=10+170/343.
\]
Both factors are lower bounds for the corresponding exponentials by
positive series. Clearing denominators gives the integer comparison
\[
\begin{split}
449233538381596942864443106714320805345633416\quad\quad\\
{}>447728960659239231243811144335360000000000000.
\end{split}
\]
For $y=x^2+\beta x+2$,
\[
 P_0(y)=\frac43x^2-\frac23\beta x-\frac43,\qquad
 \frac{P_0(y)}{x^2}\ge\frac{8024}{6075}>\frac{13}{10},
\]
where $\beta<5/3$ and $x\ge90$ were used. Also
\[
 \frac{\varepsilon(4y+\sqrt6x^3)}{x^2}
 <\frac{49}{40}+\frac2x+\frac{10}{3x^2}+\frac4{x^3}
 \le\frac{49}{40}+\frac1{45}+\frac1{2430}+\frac1{182250}<\frac54.
\]
Thus $P_\varepsilon(y)>x^2/20>0$. The positive-root argument from
Lemma~\ref{lem:scalar} proves Theorem~\ref{thm:weak}.

\subsection{Bounded difference multiplicity}
For the hypothesis of Theorem~\ref{thm:gthin}, ordered-pair expansion
and the lattice majorant give, at every $0<T\le N$,
\[
 E_f(\mu,\mu)=\frac43k+2\sum_{d\ge1}r_A(d)f(d/T)
 \le\frac43k+gT.
\]
Only the nonzero differences receive the factor $g$.
Substitute $S=gT$, $X=gN$ in the inequality of Lemma~\ref{lem:6}; every
$0<S\le X$ is allowed. It becomes~\eqref{eq:scalar}, and
Lemma~\ref{lem:scalar} proves~\eqref{eq:gthin}.
Translation to the closed interval $[0,D]$ proves the diameter version.
Empty and singleton sets are immediate; that version requires $D>0$.
For example, if $k\ge20366$, counting gives
$gD\ge k(k-1)/2\ge207376795>120^4$, so the explicit inverse is
\[
 gD>\left(\frac{\sqrt{4(k-1)+8/9}-2\sqrt2/3}{2}\right)^4.
\]

\subsection{Difference triangle sets}
Let the scope be $m$ and put $K=k+1$. For $0<T\le m$, apply the
capacity bound in Lemma~\ref{lem:6} separately to each row measure
$\mu_i=\sum_{j=0}^k\delta_{a_{ij}/T}$. Summation, using the single
shared budget for all positive within-row differences, gives
\[
 nK^2\le C(m/T)\sum_iE_f(\mu_i,\mu_i)
          \le C(m/T)(4nK/3+T).
\]
With $S=T/n$, $X=m/n$, division by $n$ gives~\eqref{eq:scalar}
for $K$, for every real $0<S\le X$. No cross-row differences are
inserted and no integrality of $X$ is needed.
Counting the $n\binom{k+1}{2}$ different positive integers yields
\[
 X\ge k(k+1)/2,\qquad
 20364\cdot20365/2=207356430<120^4<207376795=20365\cdot20366/2.
\]
Thus $k\ge20365$ allows Lemma~\ref{lem:scalar}. It gives
$k<X^{1/2}+\gamma X^{1/4}$, with $\gamma=2\sqrt2/3$.
Inverting the increasing quadratic proves $m>n\Phi(k)^4$.
For an explicit expansion put $z=\Phi(k)$, so $z^2+\gamma z=k$.
Then
\[
 z^4=k^2+\gamma^2k-\gamma(2k+\gamma^2)z,
 \qquad z\le\sqrt k-\gamma/2+\frac{\gamma^2}{8\sqrt k}.
\]
The latter follows from $\sqrt{4k+\gamma^2}\le2\sqrt k+\gamma^2/(4\sqrt k)$.
Substitution gives
\[
 z^4\ge k^2-2\gamma k^{3/2}+2\gamma^2k
       -\frac{5\gamma^3}{4}\sqrt k+\frac{\gamma^4}{2}
       -\frac{\gamma^5}{8\sqrt k}
 \ge k^2-\frac{4\sqrt2}{3}k^{3/2}+\frac{16}{9}k-2\sqrt k.
\]
For the last step use $\gamma<1$, $k\ge1$, and
$5\gamma^3/4+\gamma^5/(8k)<2$, discarding the positive constant.
This proves both statements in Theorem~\ref{thm:dts}.

\section{Product kernels}\label{sec:products}
Let $\rho$ be the pushforward of $\nu_{R/T}$ under $s\mapsto Ts$,
where $0<T\le R$. For $F_T(z)=\prod_{j=1}^d f(z_j/T)$,
the signed product measure $\rho^{\otimes d}$ has potential one on
$[0,R]^d$ and energy at most $(R/T+b+200e^{-\alpha R/T})^d$.
This follows by Fubini; each factor energy is nonnegative.
The kernel is the autocorrelation of
$T^{-d/2}\prod_j h(z_j/T)$, so signed Cauchy--Schwarz is valid.
A positive point measure $\mu$ of mass $k$ in the box therefore obeys
\begin{equation}\label{eq:product-capacity}
 k^2\le(R/T+b+200e^{-\alpha R/T})^d E_{F_T}(\mu,\mu).
\end{equation}
The certificate need not be supported in the box; its exact potential
on the box and its finite total variation are what is used.

\subsection{Manhattan diameter and the index-two lattice}
The map $(x,y)\mapsto(u,v)=(x+y,x-y)$ satisfies
$|\Delta x|+|\Delta y|=\max(|\Delta u|,|\Delta v|)$.
For a nonempty configuration translate each new coordinate separately
to lie in $[0,r]$. The transformed points may lie in either coset of
$\Lambda=\{(u,v)\in\Z^2:u\equiv v\pmod2\}$, but every difference
lies in $\Lambda$ and remains unique. Empty configurations are immediate.
Put $a=4/3$ and
\[
 S_e=\sum_{j\in\Z}f(2j/T),\qquad
 S_o=\sum_{j\in\Z}f((2j+1)/T).
\]
Monotonicity gives $S_e\le T/2+a$.
With $v_0=f(1/T)$, comparison with the preceding intervals gives
$S_o\le2v_0+T\int_{1/T}^{\infty}f\le T/2+v_0\le T/2+a$;
this is valid also for $T<1$ when the relevant samples vanish.
The complete lattice sum is $S_e^2+S_o^2$. Its zero-vector term is
$a^2$, while the point measure has $m$ such diagonal terms. Thus
\begin{equation}\label{eq:manhattan-scalar}
 m^2\le\left(r/T+b+200e^{-\alpha r/T}\right)^2
       \left(a^2m+T^2/2+2aT+a^2\right),\qquad0<T\le r.
\end{equation}
This explicitly retains the lower-order lattice terms.

For $A\ge0$, $B>0$, the positive root of $z^2-Az-B$ is at most
\begin{equation}\label{eq:quadratic-root}
 \sqrt B+A/2+A^2/(8\sqrt B).
\end{equation}
It follows from $\sqrt{B+A^2/4}\le\sqrt B+A^2/(8\sqrt B)$.
Put $x=r^{1/3}$, $c=(4/3)^{1/3}$, $t=\sqrt2c$, $T=tx^2$.
Then $1<t<2$, $c<4/3$. For $x\ge160$, the scale is valid and
\[
 \varepsilon x=200xe^{-\alpha x/t}\le200xe^{-x/8}
 \le32000e^{-20}<32000/2^{20}<1/32.
\]
Write $\beta=b+\varepsilon<1$, $D=x/t+\beta\le2x$ and
$B_0=T^2/2+2aT+a^2$.
Since $T/\sqrt2\le\sqrt{B_0}\le T/\sqrt2+\sqrt2a$,
\eqref{eq:quadratic-root} applied to~\eqref{eq:manhattan-scalar} yields
\[
 m\le D(T/\sqrt2+\sqrt2a)+a^2D^2/2
                    +a^4\sqrt2D^3/(8T).
\]
The leading terms are $x^3/\sqrt2+cx^2$, since
$bt/\sqrt2+a^2/(2t^2)=c$ (the summands are $2c/3$ and $c/3$).
The remaining terms obey
\[
\begin{array}{rcl}
\varepsilon tx^2/\sqrt2&\le&x/24,\\
a^2\beta x/t&\le&(16/9)x,\\
\sqrt2a x/t&\le&2x,\\
a^2\beta^2/2+\sqrt2a\beta&\le&26/9\le(13/720)x,\\
a^4\sqrt2D^3/(8T)&\le&(128/27)x.
\end{array}
\]
Their sum is at most $(18529/2160)x<9x$, proving
Theorem~\ref{thm:manhattan}.

\subsection{Integer boxes and a sharper lattice estimate}
For $T\ge1$ let $n_0=\lfloor T\rfloor$, $\eta=T-n_0$ and
$S(T)=\sum_{j\in\Z}f(j/T)$. Finite sums of first and third powers give
\begin{align}
 S(T)&=\frac43+\frac83n_0-\frac{2n_0(n_0+1)}T
                   +\frac{n_0^2(n_0+1)^2}{3T^3}\notag\\
 &=T+\frac1{3T}-\frac{2\eta(1-\eta)(1-2\eta)}{3T^2}
                 +\frac{\eta^2(1-\eta)^2}{3T^3}
 \le T+\frac7{16T}\le T+\frac1{2T}.\label{eq:sharp-lattice}
\end{align}
Indeed $\eta(1-\eta)\le1/4$ and
$\eta(1-\eta)|1-2\eta|\le1/8$: with $u=|1-2\eta|$, the latter
is $u(1-u^2)/4\le u(1-u)/2\le1/8$.
An integral endpoint contributes zero, so no separate rounding case is omitted.

Unique nonzero ordered vector differences in $[N]^d$ and
\eqref{eq:product-capacity} imply
\begin{equation}\label{eq:box-scalar}
 k^2\le D^d[a^dk+S(T)^d-a^d]
       \le a^dD^dk+D^d(T+1/(2T))^d,
 \quad D=N/T+b+200e^{-\alpha N/T},
\end{equation}
for $1\le T\le N$. The diagonal is $a^dk$, and the full lattice
budget includes the origin exactly once, explaining the subtraction.
No factor $2^d$ is added to an already ordered difference count.

Use the $x,c_d$ of~\eqref{eq:box-parameters}, and put
\[
 t=(a^d/b)^{1/(d+1)},\quad s=bt=(ab)^{d/(d+1)},\quad
 q=x^{-1},\quad K=N^{d/2},\quad T=tN/x.
\]
Then $1<t<3/2$, $0<s<1$, $a^d/t^d=s$ and $T\ge x$.
Under $x\ge\max(120,4d)$ we have $T\le N$, $q\le1/120$, $dq\le1/4$.
As before $\varepsilon=200e^{-\alpha x/t}$ satisfies
$\varepsilon x<24000/2^{20}<1/32$. Define
\[
 w=(b+\varepsilon)tq,\quad z=1/(2T^2),\quad v=(1+w)(1+z)-1.
\]
Then
\[
 0\le w\le sq+3q^2/64\le q+q^2,\qquad
 0\le z\le q^2/2,\qquad 0\le v\le sq+q^2\le q+q^2.
\]
For the last inequality the excess coefficient is at most
$3/64+1/2+(q+q^2)/2\le45/64<1$, already for $q\le1/4$.
Hence $dw,dv\le5/16<1/3$. The binomial estimate
$(1+u)^d\le(1-du)^{-1}$, $du<1$, gives
\[
 (1+w)^d\le3/2,\quad (1+v)^d\le3/2,\qquad
 (1+w)^d-1\le dw/(1-dw)\le2dq.
\]
Dividing~\eqref{eq:box-scalar} by $K^2$ gives a quadratic for
$U=k/K$, with linear coefficient $sq(1+w)^d$ and constant term
$(1+v)^d$; here $x^{d+1}=K$. Its root bound is
\begin{equation}\label{eq:box-root}
 U\le(1+v)^{d/2}+\frac{sq}{2}(1+w)^d
       +\frac{s^2q^2(1+w)^{2d}}{8(1+v)^{d/2}}.
\end{equation}
The final term is at most $9q^2/32$ and the middle term at most
$sq/2+dq^2$. Taylor's theorem, with $p=d/2\ge1$, bounds the
second derivative of $(1+v)^p$ by $(3/2)p(p-1)$ on the relevant
interval; for $1\le p<2$ use $(1+u)^{p-2}\le1$.
Since $v\le sq+q^2\le2q$,
\[
 (1+v)^{d/2}\le1+(d/2)sq+(d/2+3d^2/4)q^2.
\]
Consequently
\[
 U\le1+\frac{d+1}{2}sq+
       \left(\frac{3d^2}{4}+\frac{3d}{2}+\frac9{32}\right)q^2
 \le1+c_dq+2d^2q^2.
\]
The final comparison holds for all $d\ge2$ because
$5d^2/4-3d/2-9/32$ is positive at $2$ and increasing.
Multiplication by $K$ proves~\eqref{eq:boxes} in Theorem~\ref{thm:boxes}.
AM--GM on $(dbt+a^dt^{-d})/2$ explains the chosen scale and
coefficient. This optimizes the retained product inequality only.

\subsection{A two-dimensional bound for every positive integer}
For a strong Sidon $A\subseteq[N]^2$ with $k=|A|$, take an integer $T\ge1$ and
$r(z)=\#\{a\in A:z\in a+[T]^2\}$. It is supported on
$[N+T-1]^2$. Write $M=N+T-1$. Window intersections give exactly
\[
 \sum_zr(z)=kT^2,\qquad
 \sum_zr(z)^2=kT^2+
 \sum_{a\ne b}(T-|a_1-b_1|)_+(T-|a_2-b_2|)_+.
\]
Since $\sum_{j\in\Z}(T-|j|)_+=T^2$, uniqueness gives the full
nonzero-vector budget $T^4-T^2$. Thus
\begin{equation}\label{eq:alln-window}
 \sum_zr(z)^2\le T^4+(k-1)T^2,\qquad
 k^2\le M^2[1+(k-1)/T^2].
\end{equation}
The last step is finite Cauchy--Schwarz on $M^2$ positions. Empty
sets are covered as well, because $T^4-T^2\ge0$.
Put $x=N^{1/3}\ge1$, $T=\lceil x^2\rceil$ and $A_0=M^2/T^2$.
Then $x^3\le M\le x^3+x^2$, $A_0\le(x+1)^2$.
Drop the negative constant $-A_0$ in the quadratic and use
\eqref{eq:quadratic-root} to get
\begin{align*}
 k&\le M+A_0/2+A_0^2/(8M)\\
 &\le x^3+\frac32x^2+\frac98x+1+\frac3{4x}
                              +\frac1{2x^2}+\frac1{8x^3}
 \le x^3+\frac32x^2+\frac72x.
\end{align*}
The final difference, multiplied by $8x^3$, is
\[
 19x^4-8x^3-6x^2-4x-1=(x-1)(19x^3+11x^2+5x+1)\ge0.
\]
Let $c=(8/3)^{1/3}$. At $x\ge120$, \eqref{eq:boxes} with $d=2$
gives $k\le x^3+cx^2+8x$.
For $1\le x<120$, the integer comparison
$83^3=571787<576000=(8/3)60^3$ gives $c>83/60$, hence
$(3/2-c)x^2<(7/60)x^2<14x$.
The elementary bound is therefore less than
$x^3+cx^2+(35/2)x\le x^3+cx^2+18x$.
This proves~\eqref{eq:boxes-all}, including $N=1$ and the splice
point $N=120^3$.

\section{Verification status and limits}\label{sec:verification}
The transfer scout derivations were supplied by Claude (Anthropic).
GPT-6 Astra (OpenAI) examined those derivations, repaired
the combinatorial and remainder arguments where needed, and wrote the
complete explicit proofs. The cosine functional proof and admissible
perturbation were derived in this informed proof-writing process,
followed by a separate same-vendor (OpenAI) review in the same workflow.
These transfers were
not clean-room attacks. The author takes responsibility for the mathematics.

\RefereeStatus

The referee verdicts are mathematical model reviews, not independent
human referee reports. Source-formula checking of this manuscript is
separate from those reviews.

\begin{sloppypar}
Four of the bounds have been formalised and kernel-checked in Lean
v4.34.0-rc1 with Mathlib commit \texttt{de5ce8a9}, in the same project
as the earlier ordinary Sidon theorem
$|A|\le\sqrt N+(2\sqrt2/3)N^{1/4}+1$ ($N\ge120^4$):
\texttt{g\_thin\_second\_order} (Theorem~\ref{thm:gthin}),
\texttt{weakSidon\_second\_order} (Theorem~\ref{thm:weak}),
\texttt{sonar\_triangle\_bound} (the triangular bound
\eqref{eq:sonar-ramp} of Theorem~\ref{thm:sonar}), and
\texttt{difference\_triangle\_scope\_bound} and \texttt{difference\_\allowbreak triangle\_\allowbreak expanded\_\allowbreak bound} (Theorem~\ref{thm:dts}).
The statements are in the files \texttt{TransferStatement.lean} and
\texttt{DifferenceTriangle\allowbreak Statement.lean}; each theorem's reported
axioms are only \texttt{propext}, \texttt{Classical.choice} and
\texttt{Quot.sound}, enforced by \texttt{\#guard\_msgs} gates checked
on GitHub Actions (run 37076247531, commit \texttt{2f981a7}).
The public source at the checked revision is available
\href{https://github.com/chy4pro/automath/tree/2f981a7/lean/sidon30}{in the accompanying repository}.
The Lean proofs were written by GPT-6 Astra after the paper. The cosine
sonar bound \eqref{eq:sonar-cos}, the Manhattan and box bounds, the
all-$N$ square supplement and the kernel functional results are not
formalised.
\end{sloppypar}

The dependency-free JavaScript checkers supplement the proofs with
exact BigInt arithmetic. They check finite combinatorial conventions,
lattice identities, scale substitutions, and rational error envelopes.
They are available under \texttt{problems/capacity\_transfer/} in
the accompanying repository.
The recorded executions all returned PASS:
\begin{center}\small
\begin{tabular}{@{}p{0.28\textwidth}p{0.65\textwidth}@{}}
\toprule
Checker subject & Executed finite scope\\
\midrule
Sonar & 820 marginal identities, 1070 sonar prefixes, 41856 energy cases;
non-sonar negative control rejected.\\
Weak Sidon & 16384 subsets, 2048 weak Sidon sets, 16384 energy cases;
1558 repeated-difference instances.\\
Bounded multiplicity & 4096 subsets, 7521 set/parameter pairs,
60168 energy cases and 54 substitutions.\\
Difference triangles & 687 row families, 4122 energy cases,
72 substitutions including 54 nonintegral scope quotients.\\
Manhattan & 4096 subsets, 721 configurations, 3605 energy cases,
72 parity-lattice cases.\\
Boxes & 768 subsets in dimensions two and three, 341 Sidon sets,
1705 energy cases; symbolic lattice identity and 292 rational scales.\\
All-$N$ square supplement & 530 subsets, 3180 window identities,
1194 Sidon energy bounds, 1004 ceiling checks and two polynomial identities.\\
Cosine and perturbation & 16 and 29 exact rational checks, respectively,
including coefficient, Fourier-positivity, tail and denominator margins.\\
\bottomrule
\end{tabular}
\end{center}
The six principal combinatorial checkers cover 127840 energy-accounting cases.
These counts are not evidence for the capacity inequality itself:
the Manhattan and box checkers do not evaluate its certificate, and
the sonar checks at small parameters use a tail term so large that
they cannot meaningfully stress the capacity step. Their useful scope
is the exact marginal identities, lattice sums, combinatorial accounting
and final scalar algebra. The mathematical capacity proof is given in
Section~\ref{sec:capacity}. The referees additionally report numerical
certificate and finite-capacity tests; those are floating-point evidence,
not proofs, and were not independently rerun for this manuscript.
These finite checks do not prove an infinite-family statement or
replace the analytic arguments. The cosine checker does not perform
symbolic integration; the integral identities are proved in the text.
The perturbation checker checks the rational envelopes in its analytic
admissibility proof, not an assertion of optimality from a grid.

The literature comparisons use a bounded source audit dated
October 2, 2026. Its primary and secondary readings are distinguished
in Section~\ref{sec:results}; unread relevant sources are named there.
Absence of an earlier matching statement in that audit is not a
priority claim. No optimality is asserted for the displayed onsets
or remainder constants. The full horizontal functional infimum remains
open, and the product-template optimization does not exclude different
two-dimensional kernels or additional combinatorial information.

\paragraph{Licence.}
This note is dedicated to the public domain under
\href{https://creativecommons.org/publicdomain/zero/1.0/}{CC0 1.0}.

\begin{thebibliography}{CHO25}\footnotesize\raggedright\setlength{\itemsep}{0.5pt}
\bibitem[Chen26]{Chen26} H.~Chen, \emph{An explicit second-order bound for
Sidon sets and the limit of the fixed-kernel energy method}, 2026,
\href{https://doi.org/10.5281/zenodo.23103979}{Zenodo concept DOI:10.5281/zenodo.23103979}.

\bibitem[EGRT92]{EGRT92} P.~Erd\H{o}s, R.~Graham, I.~Ruzsa and H.~Taylor,
\emph{Combinatorica} \textbf{12} (1992), 39--44
(article on sonar-sequence bounds),
Theorem~4 and the comment following its proof.

\bibitem[ORTU14]{ORTU14} Osorio, Ruiz, Trujillo and Urbano,
\emph{Secuencias sonar y conjuntos de Sidon},
Revista de Ciencias (Universidad del Valle) \textbf{18}(1) (2014), 33--42,
\href{http://www.scielo.org.co/pdf/rcien/v18n1/v18n1a03.pdf}{Theorem~2.1}.

\bibitem[CK96]{CK96} W.~Chen and T.~Kl\o ve,
On generalization of sonar sequences and intermodulation-free frequencies
assignment sequences, \emph{Acta Math. Appl. Sinica (English Series)}
\textbf{12} (1996), 102--108,
\href{https://doi.org/10.1007/BF02009566}{doi:10.1007/BF02009566}.
Only the free preview of pp.~102--103 was inspected in the source audit.

\bibitem[BFR21]{BFR21} J.~Balogh, Z.~F\"uredi and S.~Roy,
An upper bound on the size of Sidon sets,
\href{https://arxiv.org/abs/2103.15850v2}{arXiv:2103.15850v2} (2021),
Sections~5--6. The weak-Sidon comparison refers specifically to this version.

\bibitem[CHO25]{CHO25} D.~Carter, Z.~Hunter and K.~O'Bryant,
On the diameter of finite Sidon sets,
\emph{Acta Math. Hungar.} \textbf{175} (2025), 108--126;
\href{https://arxiv.org/abs/2310.20032}{arXiv:2310.20032}, Section~4.
The journal abstract and preprint give different explicit
$g$-dependent estimates; they are not interchanged here.

\bibitem[CC97]{CC97} Chee and Colbourn,
\emph{IEEE Trans. Inform. Theory} \textbf{43} (1997)
(article on difference triangle sets),
\href{https://arxiv.org/abs/0712.2553}{arXiv:0712.2553}, Theorem~1
(quoting Kl\o ve, \emph{IEEE Trans. Inform. Theory} \textbf{34}
(1988), 355--361, Theorem~2).

\bibitem[BEMP10]{BEMP10} Blackburn, Etzion, Martin and Paterson,
\emph{IEEE Trans. Inform. Theory} \textbf{56} (2010), 1216--1229
(article on distinct-difference configurations);
\href{https://arxiv.org/abs/0811.3832v2}{arXiv:0811.3832v2}, Theorem~9.

\bibitem[Rob85]{Rob85} J.~P.~Robinson, Golomb rectangles,
\emph{IEEE Trans. Inform. Theory} \textbf{31} (1985), 781--787.
The square bound is cited secondarily through~\cite{Tru23};
the primary sonar discussion was not read.

\bibitem[Tru23]{Tru23} Trujillo,
\emph{Rev. Acad. Colomb. Cienc. Exactas Fis. Nat.} \textbf{47}(185)
(2023), 1024--1044 (survey), pp.~1037--1039.
The square comparisons there cite Robinson~\cite{Rob85} and
N.~Y.~Caicedo Bravo, \emph{Conjuntos de Sidon en dimensi\'on dos},
Ph.D. thesis, Universidad del Valle (2016); the thesis was not inspected.
\end{thebibliography}
\end{document}
